% Wygenerowano z ../00_wprowadzenie.md skryptem _md_do_latex.py (można edytować ręcznie).
\section{Jak pracujemy}\label{jak-pracujemy}

Pracujecie w zespołach 2--3-osobowych. Każdy zespół ma \textbf{jednego Jetbota}, \textbf{stanowisko komputerowe} i \textbf{ten zeszyt ćwiczeń}. Wyniki wpisujecie bezpośrednio do tabel w instrukcjach. Przydadzą się w podsumowaniu.

Podział pracy ustalacie sami. Pilnujcie tylko dwóch zasad:

\begin{itemize}
\tightlist
\item
  \textbf{przy klawiaturze zmieniajcie się co ćwiczenie}, żeby każdy coś uruchomił i zmienił w notatniku,
\item
  \textbf{gdy robot jedzie, jedna osoba stoi przy torze} (w instrukcjach nazywamy ją \textbf{pilotem}). Ustawia robota na starcie, a w razie zagrożenia łapie go albo woła „STOP!''.
\end{itemize}

\subsection{Ćwiczenia}\label{ux107wiczenia}

\par\noindent\begin{minipage}{\linewidth}
\begin{tabular}{@{}
  >{\raggedright\arraybackslash}p{(\linewidth - 6\tabcolsep) * \real{0.1053}}
  >{\raggedright\arraybackslash}p{(\linewidth - 6\tabcolsep) * \real{0.5684}}
  >{\raggedright\arraybackslash}p{(\linewidth - 6\tabcolsep) * \real{0.2211}}
  >{\raggedright\arraybackslash}p{(\linewidth - 6\tabcolsep) * \real{0.1053}}@{}}
\toprule\noalign{}
\begin{minipage}[b]{\linewidth}\raggedright
Nr
\end{minipage} & \begin{minipage}[b]{\linewidth}\raggedright
Ćwiczenie
\end{minipage} & \begin{minipage}[b]{\linewidth}\raggedright
Notatnik
\end{minipage} & \begin{minipage}[b]{\linewidth}\raggedright
Czas
\end{minipage} \\
\midrule\noalign{}
0 & Pierwsze kroki z notatnikiem i robotem & \texttt{00\_start.ipynb} & 5 min \\
1 & LiDAR: jak robot mierzy odległość & \texttt{01\_lidar.ipynb} & 12 min \\
2 & Kamera: jak zapisany jest obraz & \texttt{02\_kamera.ipynb} & 8 min \\
3 & Odometria: dlaczego robot się gubi & \texttt{03\_odometria.ipynb} & 15 min \\
4 & SLAM: budujemy mapę & \texttt{04\_slam.ipynb} & 40 min \\
& \emph{przerwa} & & \\
5 & Neuron: hamuj czy jedź? & \texttt{05\_perceptron.ipynb} & 10 min \\
6 & Gotowa sieć neuronowa szuka obiektów: jak zapytać sieć & \texttt{06\_detektor.ipynb} & 20 min \\
7 & Autonomiczna jazda na mapie i poszukiwanie kaczki & \texttt{07\_nawigacja.ipynb} & 35 min \\
8 & Karta wyników zespołu & (papier) & 20 min \\
\bottomrule\noalign{}
\end{tabular}
\end{minipage}\par\medskip

\section{Ćwiczenie 0: pierwsze kroki (5 min)}\label{ux107wiczenie-0-pierwsze-kroki-5-min}

\textbf{Cel:} nauczyć się uruchamiać komórki notatnika i sprawdzić, czy robot jest połączony.

\subsection{Jak działa notatnik Jupyter}\label{jak-dziaux142a-notatnik-jupyter}

\begin{itemize}
\tightlist
\item
  Notatnik składa się z \textbf{komórek}. Komórki z tekstem to opis, a komórki z kodem to polecenia dla komputera.
\item
  Komórkę uruchamiasz, gdy klikniesz ją i naciśniesz \textbf{Shift + Enter} (albo przycisk ▶).
\item
  Gdy komórka pracuje, obok niej widać \texttt{{[}*{]}}. Po zakończeniu pojawia się numer, np. \texttt{{[}3{]}}.
\item
  Komórki uruchamiaj \textbf{po kolei, od góry}.
\item
  Wolno zmieniać tylko wartości w liniach oznaczonych \textbf{\texttt{\#\ ZMIEŃ\ TUTAJ}}. Resztę kodu można czytać, ale nie trzeba go rozumieć.
\item
  Gdy coś „zawiśnie'', wybierz z menu \emph{Kernel → Restart} i uruchom komórki od początku.
\item
  \textbf{Ostatnia komórka każdego notatnika to \texttt{koniec()}.} Uruchom ją, zanim otworzysz następny notatnik: zatrzymuje robota i odłącza notatnik. Robot działa dalej, a następny notatnik połączy się z nim w pierwszej komórce. Niczego nie trzeba restartować ani zamykać.
\end{itemize}

\subsection{Kroki}\label{kroki}

\begin{enumerate}
\def\labelenumi{\arabic{enumi}.}
\tightlist
\item
  Otwórz \texttt{00\_start.ipynb}. Uruchom pierwszą komórkę (połączenie z robotem). Ma pojawić się komunikat \emph{„Połączono z robotem''}. \textbf{Sprawdź, czy numer w komunikacie, np. \texttt{{[}robot 3{]}}, zgadza się z naklejką na Waszym robocie.} Jeśli nie, zawołaj prowadzącego.
\item
  Uruchom komórkę \texttt{status()}. Przy napędzie, LiDAR-ze i kamerze powinien być znak \ding{51}.
\item
  W komórce z krótkim przejazdem zmień czas jazdy (\texttt{czas\ =\ 1.0\ \ \#\ ZMIEŃ\ TUTAJ}) na \texttt{2.0} i uruchom ją. \textbf{Pilot} trzyma rękę przy robocie.
\item
  Uruchom komórkę z \texttt{stop()}. Zapamiętaj ją, bo to \textbf{hamulec awaryjny}, który działa w każdym notatniku.
\end{enumerate}

\begin{quote}
\textbf{Bezpieczeństwo.} Robot jest mały, ale potrafi spaść ze stołu albo wjechać pod nogi. Gdy robot jedzie, pilot jest zawsze przy torze. W razie problemu: komórka \texttt{stop()} albo podnieś robota.
\end{quote}
